\chapter{总结与展望}
\enchapter{Conclusion and Prospect}

\section{全文总结}
\ensection{Summary of Full-text Work}

本文围绕碳纳米管阵列（CNTA）的制备工艺优化，结合领域知识增强与机器学习预测技术，完成了从知识库构建、形貌特征提取到集成化平台实现的一系列研究工作。主要总结如下：

首先，针对 CNT 领域知识来源分散且利用率低的问题，构建了领域知识库，并设计了任务感知的检索增强机制，使多源异构知识能够被统一组织和检索。

其次，针对 SEM 图片形貌量化难的问题，开发了自动化特征计算算法，为形貌预测整理出结构化数据集。

最后，构建了基于工艺参数的形貌预测模型，并将其集成到综合智能分析平台中，形成从实验数据到分析和预测的连续流程。

\section{研究工作总结与创新点归纳}
\ensection{Summary of Innovations}

1. **知识增强的分析框架**：不同于纯数据驱动的材料预测，本文引入 RAG 技术，将预测结果与相关机理和文献证据联系起来。

2. **面向科研流程的方法集成**：通过集成化平台建设，将 CNT 研究中的知识检索、图像表征与参数预测连接起来。

\section{后续研究展望}
\ensection{Future Research Prospects}

尽管本文已完成上述工作，但仍有以下方向值得继续研究：
1. 引入多模态学习技术，实现实验曲线、能谱等多维度数据的融合建模。
2. 探索预测模型在在线工艺控制中的原位应用。
3. 进一步增强大模型对复杂材料机理的推理一致性。
